\documentclass[10pt,letterpaper]{article}
\usepackage{cornell-notes}
\usepackage{mathtools}

\title{Chapter 19: Integral Equations and Inverse Theory}
\author{}
\date{}

\setDocTitle{Chapter 19: Integral Equations and Inverse Theory}
\setDocSubtitle{Numerical Methods}
\setDocVersion{v1.0}
\setDocStatus{Study Companion}
\setDocOwner{Numerical Methods Cornell Notes}
\setDocAuthor{}
\setDocDate{\today}
\setDocSummary{Cornell-format study and review notes for Chapter 19: Integral Equations and Inverse Theory.}

\setCornellCollection{Numerical Methods Cornell Notes}
\setCornellUnitType{Chapter}
\setCornellUnitNumber{19}
\setCornellUnitTitle{Integral Equations and Inverse Theory}
\setCornellCompanionLabel{Study and review companion}

\hypersetup{pdftitle={Chapter 19: Integral Equations and Inverse Theory},pdfsubject={Cornell notes},pdfauthor={}}

\begin{document}
\maketitle

\section*{Chapter Overview}
This chapter organizes the mathematical and computational ideas behind integral equations and inverse theory. The notes preserve the numbered section sequence while emphasizing method selection, explicit assumptions, numerical reliability, cost, and verification.

\section*{Learning Objectives}
Explain the governing problem class; compare Fredholm Equations of the Second Kind, Volterra Equations, Integral Equations with Singular Kernels, and Inverse Problems and the Use of A Priori Information; design defensible stopping and verification checks; distinguish problem conditioning from algorithmic stability.

\section*{Prerequisites}
Integration, linear algebra, singular values, probability, and regularization concepts.

\section*{Operational Workflow}
Identify equation kind and kernel; discretize with compatible quadrature; inspect singular values or conditioning; introduce prior structure explicitly; select regularization; validate resolution, residual, and sensitivity.

\subsection*{Formula and notation anchor}
\[
\widehat{x}_\lambda=\arg\min_x\left(\|Ax-b\|_2^2+\lambda^2\|Lx\|_2^2\right)
\]
The parameter $\lambda$ balances fidelity against stability; $L$ defines which solution features are penalized.

\begin{warnbox}{Numerical warning}
Noise amplification, discretization-dependent priors, overfitting the data, oversmoothing, invalid positivity assumptions, singular kernels, and reporting one reconstruction without resolution information.
\end{warnbox}

\section{19.0 Introduction}
\begin{cornell}
\noterow{What is the central idea?}{Integral equations and inverse problems transform unknown functions through kernels; smoothing often makes inversion ill-conditioned and dependent on prior information.}
\noterow{How should it be used?}{For Introduction, begin from explicit inputs, outputs, preconditions, and representation choices.  Identify equation kind and kernel; discretize with compatible quadrature; inspect singular values or conditioning; introduce prior structure explicitly; select regularization; validate resolution, residual, and sensitivity.}
\noterow{Cost and reliability}{Analyze arithmetic work, storage, data movement, and convergence separately.  Relevant failure pressures include noise amplification, discretization-dependent priors, overfitting the data, oversmoothing, invalid positivity assumptions, singular kernels, and reporting one reconstruction without resolution information.  Use scaled tolerances and report both success criteria and diagnostic evidence.}
\noterow{Implementation checklist}{\begin{itemize}
\item Validate dimensions, domains, ordering, and structural assumptions.
\item Guard small divisors, invalid states, overflow, underflow, and nonfinite values.
\item Make mutation, workspace, indexing, and termination behavior explicit.
\item Test a simple exact case, a difficult valid case, a boundary case, and an expected failure.
\end{itemize}}
\end{cornell}
\begin{summarybox}{Section summary}
Integral equations and inverse problems transform unknown functions through kernels; smoothing often makes inversion ill-conditioned and dependent on prior information.  Select this technique only when its assumptions match the data, and accept a result only after an independent residual, error estimate, invariant, or refinement check.
\end{summarybox}

\section{19.1 Fredholm Equations of the Second Kind}
\begin{cornell}
\noterow{What is the central idea?}{A second-kind equation contains the unknown both directly and under an integral, often yielding a better-conditioned discretized linear system.}
\noterow{How should it be used?}{For Fredholm Equations of the Second Kind, begin from explicit inputs, outputs, preconditions, and representation choices.  Identify equation kind and kernel; discretize with compatible quadrature; inspect singular values or conditioning; introduce prior structure explicitly; select regularization; validate resolution, residual, and sensitivity.}
\noterow{Quantitative anchor}{\[
f(x)=g(x)+\lambda\int_a^b K(x,t)f(t)\,dt
\]
State the dimensions, scaling, and validity assumptions before using this relation.  Check the resulting residual or invariant rather than trusting an iteration count alone.}
\noterow{Cost and reliability}{Analyze arithmetic work, storage, data movement, and convergence separately.  Relevant failure pressures include noise amplification, discretization-dependent priors, overfitting the data, oversmoothing, invalid positivity assumptions, singular kernels, and reporting one reconstruction without resolution information.  Use scaled tolerances and report both success criteria and diagnostic evidence.}
\noterow{Implementation checklist}{\begin{itemize}
\item Validate dimensions, domains, ordering, and structural assumptions.
\item Guard small divisors, invalid states, overflow, underflow, and nonfinite values.
\item Make mutation, workspace, indexing, and termination behavior explicit.
\item Test a simple exact case, a difficult valid case, a boundary case, and an expected failure.
\end{itemize}}
\end{cornell}
\begin{summarybox}{Section summary}
A second-kind equation contains the unknown both directly and under an integral, often yielding a better-conditioned discretized linear system.  Select this technique only when its assumptions match the data, and accept a result only after an independent residual, error estimate, invariant, or refinement check.
\end{summarybox}

\section{19.2 Volterra Equations}
\begin{cornell}
\noterow{What is the central idea?}{A variable upper limit gives Volterra equations a causal or triangular structure that supports forward progression.}
\noterow{How should it be used?}{For Volterra Equations, begin from explicit inputs, outputs, preconditions, and representation choices.  Identify equation kind and kernel; discretize with compatible quadrature; inspect singular values or conditioning; introduce prior structure explicitly; select regularization; validate resolution, residual, and sensitivity.}
\noterow{Quantitative anchor}{\[
f(x)=g(x)+\int_a^xK(x,t)f(t)\,dt
\]
State the dimensions, scaling, and validity assumptions before using this relation.  Check the resulting residual or invariant rather than trusting an iteration count alone.}
\noterow{Cost and reliability}{Analyze arithmetic work, storage, data movement, and convergence separately.  Relevant failure pressures include noise amplification, discretization-dependent priors, overfitting the data, oversmoothing, invalid positivity assumptions, singular kernels, and reporting one reconstruction without resolution information.  Use scaled tolerances and report both success criteria and diagnostic evidence.}
\noterow{Implementation checklist}{\begin{itemize}
\item Validate dimensions, domains, ordering, and structural assumptions.
\item Guard small divisors, invalid states, overflow, underflow, and nonfinite values.
\item Make mutation, workspace, indexing, and termination behavior explicit.
\item Test a simple exact case, a difficult valid case, a boundary case, and an expected failure.
\end{itemize}}
\end{cornell}
\begin{summarybox}{Section summary}
A variable upper limit gives Volterra equations a causal or triangular structure that supports forward progression.  Select this technique only when its assumptions match the data, and accept a result only after an independent residual, error estimate, invariant, or refinement check.
\end{summarybox}

\section{19.3 Integral Equations with Singular Kernels}
\begin{cornell}
\noterow{What is the central idea?}{Known kernel singularities should be isolated or integrated analytically so quadrature acts on a regular remainder.}
\noterow{How should it be used?}{For Integral Equations with Singular Kernels, begin from explicit inputs, outputs, preconditions, and representation choices.  Identify equation kind and kernel; discretize with compatible quadrature; inspect singular values or conditioning; introduce prior structure explicitly; select regularization; validate resolution, residual, and sensitivity.}
\noterow{Quantitative lens}{Identify the governing equation, recurrence, probability law, or geometric predicate for Integral Equations with Singular Kernels.  Define every variable and scale; then derive a residual, error estimate, or invariant that can be checked independently.}
\noterow{Cost and reliability}{Analyze arithmetic work, storage, data movement, and convergence separately.  Relevant failure pressures include noise amplification, discretization-dependent priors, overfitting the data, oversmoothing, invalid positivity assumptions, singular kernels, and reporting one reconstruction without resolution information.  Use scaled tolerances and report both success criteria and diagnostic evidence.}
\noterow{Implementation checklist}{\begin{itemize}
\item Validate dimensions, domains, ordering, and structural assumptions.
\item Guard small divisors, invalid states, overflow, underflow, and nonfinite values.
\item Make mutation, workspace, indexing, and termination behavior explicit.
\item Test a simple exact case, a difficult valid case, a boundary case, and an expected failure.
\end{itemize}}
\end{cornell}
\begin{summarybox}{Section summary}
Known kernel singularities should be isolated or integrated analytically so quadrature acts on a regular remainder.  Select this technique only when its assumptions match the data, and accept a result only after an independent residual, error estimate, invariant, or refinement check.
\end{summarybox}

\section{19.4 Inverse Problems and the Use of A Priori Information}
\begin{cornell}
\noterow{What is the central idea?}{Prior smoothness, scale, positivity, or structure restricts otherwise unstable solutions but introduces an explicit bias-variance tradeoff.}
\noterow{How should it be used?}{For Inverse Problems and the Use of A Priori Information, begin from explicit inputs, outputs, preconditions, and representation choices.  Identify equation kind and kernel; discretize with compatible quadrature; inspect singular values or conditioning; introduce prior structure explicitly; select regularization; validate resolution, residual, and sensitivity.}
\noterow{Quantitative lens}{Identify the governing equation, recurrence, probability law, or geometric predicate for Inverse Problems and the Use of A Priori Information.  Define every variable and scale; then derive a residual, error estimate, or invariant that can be checked independently.}
\noterow{Cost and reliability}{Analyze arithmetic work, storage, data movement, and convergence separately.  Relevant failure pressures include noise amplification, discretization-dependent priors, overfitting the data, oversmoothing, invalid positivity assumptions, singular kernels, and reporting one reconstruction without resolution information.  Use scaled tolerances and report both success criteria and diagnostic evidence.}
\noterow{Implementation checklist}{\begin{itemize}
\item Validate dimensions, domains, ordering, and structural assumptions.
\item Guard small divisors, invalid states, overflow, underflow, and nonfinite values.
\item Make mutation, workspace, indexing, and termination behavior explicit.
\item Test a simple exact case, a difficult valid case, a boundary case, and an expected failure.
\end{itemize}}
\end{cornell}
\begin{summarybox}{Section summary}
Prior smoothness, scale, positivity, or structure restricts otherwise unstable solutions but introduces an explicit bias-variance tradeoff.  Select this technique only when its assumptions match the data, and accept a result only after an independent residual, error estimate, invariant, or refinement check.
\end{summarybox}

\section{19.5 Linear Regularization Methods}
\begin{cornell}
\noterow{What is the central idea?}{Regularization minimizes data misfit plus a stabilizing penalty, with the regularization strength controlling fit versus sensitivity.}
\noterow{How should it be used?}{For Linear Regularization Methods, begin from explicit inputs, outputs, preconditions, and representation choices.  Identify equation kind and kernel; discretize with compatible quadrature; inspect singular values or conditioning; introduce prior structure explicitly; select regularization; validate resolution, residual, and sensitivity.}
\noterow{Quantitative anchor}{\[
\min_x\|Ax-b\|_2^2+\lambda^2\|Lx\|_2^2
\]
State the dimensions, scaling, and validity assumptions before using this relation.  Check the resulting residual or invariant rather than trusting an iteration count alone.}
\noterow{Cost and reliability}{Analyze arithmetic work, storage, data movement, and convergence separately.  Relevant failure pressures include noise amplification, discretization-dependent priors, overfitting the data, oversmoothing, invalid positivity assumptions, singular kernels, and reporting one reconstruction without resolution information.  Use scaled tolerances and report both success criteria and diagnostic evidence.}
\noterow{Implementation checklist}{\begin{itemize}
\item Validate dimensions, domains, ordering, and structural assumptions.
\item Guard small divisors, invalid states, overflow, underflow, and nonfinite values.
\item Make mutation, workspace, indexing, and termination behavior explicit.
\item Test a simple exact case, a difficult valid case, a boundary case, and an expected failure.
\end{itemize}}
\end{cornell}
\begin{summarybox}{Section summary}
Regularization minimizes data misfit plus a stabilizing penalty, with the regularization strength controlling fit versus sensitivity.  Select this technique only when its assumptions match the data, and accept a result only after an independent residual, error estimate, invariant, or refinement check.
\end{summarybox}

\section{19.6 Backus-Gilbert Method}
\begin{cornell}
\noterow{What is the central idea?}{Backus-Gilbert inversion constructs localized averaging kernels and quantifies the tradeoff between resolution and estimator variance.}
\noterow{How should it be used?}{For Backus-Gilbert Method, begin from explicit inputs, outputs, preconditions, and representation choices.  Identify equation kind and kernel; discretize with compatible quadrature; inspect singular values or conditioning; introduce prior structure explicitly; select regularization; validate resolution, residual, and sensitivity.}
\noterow{Quantitative lens}{Identify the governing equation, recurrence, probability law, or geometric predicate for Backus-Gilbert Method.  Define every variable and scale; then derive a residual, error estimate, or invariant that can be checked independently.}
\noterow{Cost and reliability}{Analyze arithmetic work, storage, data movement, and convergence separately.  Relevant failure pressures include noise amplification, discretization-dependent priors, overfitting the data, oversmoothing, invalid positivity assumptions, singular kernels, and reporting one reconstruction without resolution information.  Use scaled tolerances and report both success criteria and diagnostic evidence.}
\noterow{Implementation checklist}{\begin{itemize}
\item Validate dimensions, domains, ordering, and structural assumptions.
\item Guard small divisors, invalid states, overflow, underflow, and nonfinite values.
\item Make mutation, workspace, indexing, and termination behavior explicit.
\item Test a simple exact case, a difficult valid case, a boundary case, and an expected failure.
\end{itemize}}
\end{cornell}
\begin{summarybox}{Section summary}
Backus-Gilbert inversion constructs localized averaging kernels and quantifies the tradeoff between resolution and estimator variance.  Select this technique only when its assumptions match the data, and accept a result only after an independent residual, error estimate, invariant, or refinement check.
\end{summarybox}

\section{19.7 Maximum Entropy Image Restoration}
\begin{cornell}
\noterow{What is the central idea?}{Maximum-entropy restoration enforces positivity and selects a reconstruction by balancing data agreement with an entropy-based prior.}
\noterow{How should it be used?}{For Maximum Entropy Image Restoration, begin from explicit inputs, outputs, preconditions, and representation choices.  Identify equation kind and kernel; discretize with compatible quadrature; inspect singular values or conditioning; introduce prior structure explicitly; select regularization; validate resolution, residual, and sensitivity.}
\noterow{Quantitative lens}{Identify the governing equation, recurrence, probability law, or geometric predicate for Maximum Entropy Image Restoration.  Define every variable and scale; then derive a residual, error estimate, or invariant that can be checked independently.}
\noterow{Cost and reliability}{Analyze arithmetic work, storage, data movement, and convergence separately.  Relevant failure pressures include noise amplification, discretization-dependent priors, overfitting the data, oversmoothing, invalid positivity assumptions, singular kernels, and reporting one reconstruction without resolution information.  Use scaled tolerances and report both success criteria and diagnostic evidence.}
\noterow{Implementation checklist}{\begin{itemize}
\item Validate dimensions, domains, ordering, and structural assumptions.
\item Guard small divisors, invalid states, overflow, underflow, and nonfinite values.
\item Make mutation, workspace, indexing, and termination behavior explicit.
\item Test a simple exact case, a difficult valid case, a boundary case, and an expected failure.
\end{itemize}}
\end{cornell}
\begin{summarybox}{Section summary}
Maximum-entropy restoration enforces positivity and selects a reconstruction by balancing data agreement with an entropy-based prior.  Select this technique only when its assumptions match the data, and accept a result only after an independent residual, error estimate, invariant, or refinement check.
\end{summarybox}

\section*{Method comparison}
\begin{center}
\small
\begin{tabularx}{\linewidth}{>{\bfseries\raggedright\arraybackslash}p{0.22\linewidth}XX}
\toprule
Method or lens & Best use & Main caution \\
\midrule
Second-kind equation & Unknown appears inside and outside integral & Often better conditioned \\
Tikhonov regularization & Linear noisy inverse problem & Explicit bias-stability tradeoff \\
Maximum entropy & Positive image-like unknown & Nonlinear prior and optimization \\
\bottomrule
\end{tabularx}
\end{center}

\section*{Worked example}
\begin{exambox}{Independent worked example}
For scalar $A=0.01$, noisy data $b=0.0101$ imply the unregularized estimate $x=1.01$.  With $L=1$ and $\lambda=0.01$, $x_\lambda=Ab/(A^2+\lambda^2)=0.505$.  Stability improves but bias increases; the example makes the tradeoff visible rather than declaring one estimate universally better.
\end{exambox}

\section*{Chapter synthesis}
\begin{summarybox}{Chapter synthesis}
The chapter's methods should be viewed as a decision system rather than a list of routines.  Begin with the mathematical problem and its conditioning, exploit structure, select an algorithm with compatible stability and cost, and verify the computed answer with evidence that is independent of the internal iteration.  The recurring implementation discipline is: identify equation kind and kernel; discretize with compatible quadrature; inspect singular values or conditioning; introduce prior structure explicitly; select regularization; validate resolution, residual, and sensitivity.
\end{summarybox}

\subsection*{Common mistakes and numerical hazards}
\begin{warnbox}{Common mistakes and numerical hazards}
Noise amplification, discretization-dependent priors, overfitting the data, oversmoothing, invalid positivity assumptions, singular kernels, and reporting one reconstruction without resolution information.  A second recurring mistake is reporting only a computed value: always retain tolerances, iteration counts, residuals, refinement behavior, and relevant structural checks.
\end{warnbox}

\subsection*{Retrieval practice}
\textbf{Questions}
\begin{enumerate}
\item What problem does Fredholm Equations of the Second Kind address, and which assumption is easiest to violate?
\item What problem does Volterra Equations address, and which assumption is easiest to violate?
\item What problem does Integral Equations with Singular Kernels address, and which assumption is easiest to violate?
\item What problem does Backus-Gilbert Method address, and which assumption is easiest to violate?
\item What problem does Maximum Entropy Image Restoration address, and which assumption is easiest to violate?
\item How do conditioning, stability, and the final residual play different roles in this chapter?
\end{enumerate}

\textbf{Brief answers}
\begin{enumerate}
\item A second-kind equation contains the unknown both directly and under an integral, often yielding a better-conditioned discretized linear system. Recheck the section's domain, structure, scaling, and failure diagnostics.
\item A variable upper limit gives Volterra equations a causal or triangular structure that supports forward progression. Recheck the section's domain, structure, scaling, and failure diagnostics.
\item Known kernel singularities should be isolated or integrated analytically so quadrature acts on a regular remainder. Recheck the section's domain, structure, scaling, and failure diagnostics.
\item Backus-Gilbert inversion constructs localized averaging kernels and quantifies the tradeoff between resolution and estimator variance. Recheck the section's domain, structure, scaling, and failure diagnostics.
\item Maximum-entropy restoration enforces positivity and selects a reconstruction by balancing data agreement with an entropy-based prior. Recheck the section's domain, structure, scaling, and failure diagnostics.
\item Conditioning describes sensitivity of the mathematical problem, stability describes error propagation by the algorithm, and the residual measures satisfaction of the computed equations; none is a universal substitute for the others.
\end{enumerate}

\subsection*{Mastery checklist}
\begin{itemize}
\item[$\square$] I can explain and select Fredholm Equations of the Second Kind without relying on a memorized code listing.
\item[$\square$] I can explain and select Volterra Equations without relying on a memorized code listing.
\item[$\square$] I can explain and select Integral Equations with Singular Kernels without relying on a memorized code listing.
\item[$\square$] I can explain and select Inverse Problems and the Use of A Priori Information without relying on a memorized code listing.
\item[$\square$] I can explain and select Linear Regularization Methods without relying on a memorized code listing.
\item[$\square$] I can state a scaled stopping rule and an independent verification check.
\item[$\square$] I can identify at least one conditioning risk and one algorithmic stability risk.
\item[$\square$] I can estimate time, storage, and data-movement costs at the appropriate level.
\end{itemize}

\end{document}

